% GENERATED FILE. Edit canonical lessons, then run npm run generate:books.
% canonical-source-hash: fnv1a64:50936579e8d06c8d
% canonical-lessons: GU-C524-sangrah, GU-C524-bhamardo, GU-C524-kodi, GU-C524-santakukdi, GU-C524-koydo, GU-R524-review-517-524

\chapter{Games and Collections}
\label{ch:gu-games-and-collections}
\hlchaptermodality{524}

\begin{chapteropening}
\textbf{By the end of this chapter:} \emph{I can talk about collecting things and children's games.}

Complete the last lesson of chapter 524: I can talk about collecting things and children's games.
\end{chapteropening}

\section[\texorpdfstring{\gu{સંગ્રહ} (sangrah)}{સંગ્રહ (sangrah)}]{\gu{સંગ્રહ} (sangrah) — a collection}

\label{lesson:GU-C524-sangrah}



\begin{quote}
Before the new one: say the Gujarati for to be to one's taste, then the Gujarati for to dunk.
\end{quote}

\subsection*{You'll want to know: \gu{સંગ્રહ}}
\textbf{\gu{સંગ્રહ}} — \emph{sangrah} — \textquotedblleft{}a collection\textquotedblright{}.

\begin{quote}
\gu{અરુણ} \gu{પાસે} \gu{જૂના} \gu{સિક્કાઓનો} \gu{મોટો} \gu{સંગ્રહ} \gu{છે}.
\end{quote}

\begin{quote}
\emph{Aruṇ pāse jūnā sikkāono moṭo sangrah che.}
\end{quote}

\begin{quote}
Arun has a big collection of old coins.
\end{quote}

A masculine noun.

\begin{cousinweb}
\textbf{\gu{સંગ્રહ}} (\emph{sangrah}) is a tatsama, taken from Sanskrit \emph{saṃgraha}, \textquotedblleft{}holding together, gathering\textquotedblright{} (Monier-Williams).
\end{cousinweb}

\begin{grammarlens}[title={What you've built}]
One more word for describing what you lived and felt.
\end{grammarlens}

\subsection*{Your turn}
\begin{itemize}
\raggedright
  \item \emph{Say it:} \emph{sangrah}
  \item \emph{Say it:} \emph{sangrah}, once more
  \item \emph{Say it:} \emph{Aruṇ pāse jūnā sikkāono moṭo sangrah che.} — Arun has a big collection of old coins.
  \item \emph{Recall:} say the Gujarati for to be to one's taste, then the Gujarati for to dunk, then say \emph{sangrah} again
\end{itemize}

\begin{tcolorbox}[breakable,colback=teal!4,colframe=teal!35!black,title={Before you move on}]
What does \gu{સંગ્રહ} mean? (\textquotedblleft{}A collection\textquotedblright{}.) Say it once more.
\end{tcolorbox}
\section[\texorpdfstring{\gu{ભમરડો} (bhamarḍo)}{ભમરડો (bhamarḍo)}]{\gu{ભમરડો} (bhamarḍo) — a spinning top}

\label{lesson:GU-C524-bhamardo}



\begin{quote}
Before the new one: say the Gujarati for to dunk, then the Gujarati for a collection.
\end{quote}

\subsection*{You'll want to know: \gu{ભમરડો}}
\textbf{\gu{ભમરડો}} — \emph{bhamarḍo} — \textquotedblleft{}a spinning top\textquotedblright{}.

\begin{quote}
\gu{મારો} \gu{નાનો} \gu{ભાઈ} \gu{ભમરડા} \gu{સાથે} \gu{રમે} \gu{છે}.
\end{quote}

\begin{quote}
\emph{māro nāno bhāī bhamarḍā sāthe rame che.}
\end{quote}

\begin{quote}
My little brother plays with a spinning top.
\end{quote}

A masculine noun.

\begin{cousinweb}
\textbf{\gu{ભમરડો}} (\emph{bhamarḍo}) is filed by CDIAL under Sanskrit \emph{bhramara}, \textquotedblleft{}whirling; a whirlpool\textquotedblright{}, and CDIAL lists the Gujarati word for a spinning top (CDIAL 9650).
\end{cousinweb}

\begin{grammarlens}[title={What you've built}]
One more word for describing what you lived and felt.
\end{grammarlens}

\subsection*{Your turn}
\begin{itemize}
\raggedright
  \item \emph{Say it:} \emph{bhamarḍo}
  \item \emph{Say it:} \emph{bhamarḍo}, once more
  \item \emph{Say it:} \emph{māro nāno bhāī bhamarḍā sāthe rame che.} — My little brother plays with a spinning top.
  \item \emph{Recall:} say the Gujarati for to dunk, then the Gujarati for a collection, then say \emph{bhamarḍo} again
\end{itemize}

\begin{tcolorbox}[breakable,colback=teal!4,colframe=teal!35!black,title={Before you move on}]
What does \gu{ભમરડો} mean? (\textquotedblleft{}A spinning top\textquotedblright{}.) Say it once more.
\end{tcolorbox}
\section[\texorpdfstring{\gu{કોડી} (koḍī)}{કોડી (koḍī)}]{\gu{કોડી} (koḍī) — a cowrie shell}

\label{lesson:GU-C524-kodi}



\begin{quote}
Before the new one: say the Gujarati for a collection, then the Gujarati for a spinning top.
\end{quote}

\subsection*{You'll want to know: \gu{કોડી}}
\textbf{\gu{કોડી}} — \emph{koḍī} — \textquotedblleft{}a cowrie shell\textquotedblright{}.

\begin{quote}
\gu{દાદી} \gu{કોડીઓથી} \gu{એક} \gu{જૂની} \gu{રમત} \gu{રમતી} \gu{હતી}.
\end{quote}

\begin{quote}
\emph{dādī koḍīothī ek jūnī ramat ramtī hatī.}
\end{quote}

\begin{quote}
Grandmother used to play an old game with cowrie shells.
\end{quote}

A feminine noun.

\begin{cousinweb}
\textbf{\gu{કોડી}} (\emph{koḍī}) goes back to Sanskrit \emph{kaparda}, \textquotedblleft{}a small shell\textquotedblright{}, which CDIAL notes may itself come from Dravidian; CDIAL lists the Gujarati word for a cowrie (CDIAL 2740). Cowries were once counters in games, and small change.
\end{cousinweb}

\begin{grammarlens}[title={What you've built}]
One more word for describing what you lived and felt.
\end{grammarlens}

\subsection*{Your turn}
\begin{itemize}
\raggedright
  \item \emph{Say it:} \emph{koḍī}
  \item \emph{Say it:} \emph{koḍī}, once more
  \item \emph{Say it:} \emph{dādī koḍīothī ek jūnī ramat ramtī hatī.} — Grandmother used to play an old game with cowrie shells.
  \item \emph{Recall:} say the Gujarati for a collection, then the Gujarati for a spinning top, then say \emph{koḍī} again
\end{itemize}

\begin{tcolorbox}[breakable,colback=teal!4,colframe=teal!35!black,title={Before you move on}]
What does \gu{કોડી} mean? (\textquotedblleft{}A cowrie shell\textquotedblright{}.) Say it once more.
\end{tcolorbox}
\section[\texorpdfstring{\gu{સંતાકૂકડી} (santākūkḍī)}{સંતાકૂકડી (santākūkḍī)}]{\gu{સંતાકૂકડી} (santākūkḍī) — hide-and-seek}

\label{lesson:GU-C524-santakukdi}



\begin{quote}
Before the new one: say the Gujarati for a spinning top, then the Gujarati for a cowrie.
\end{quote}

\subsection*{You'll want to know: \gu{સંતાકૂકડી}}
\textbf{\gu{સંતાકૂકડી}} — \emph{santākūkḍī} — \textquotedblleft{}hide-and-seek\textquotedblright{}.

\begin{quote}
\gu{બાળકો} \gu{બગીચામાં} \gu{સંતાકૂકડી} \gu{રમતાં} \gu{હતાં}.
\end{quote}

\begin{quote}
\emph{bāḷko bagīcāmā̃ santākūkḍī ramtā̃ hatā̃.}
\end{quote}

\begin{quote}
The children were playing hide-and-seek in the garden.
\end{quote}

A feminine noun.

\begin{cousinweb}
\textbf{\gu{સંતાકૂકડી}} (\emph{santākūkḍī}) starts from \emph{santāvũ}, \textquotedblleft{}to hide\textquotedblright{}, which CDIAL files under Sanskrit \emph{sáṁtāyatē} (CDIAL 12887); the second half is not explained in the dictionaries we use, so its history is not settled.
\end{cousinweb}

\begin{grammarlens}[title={What you've built}]
One more word for describing what you lived and felt.
\end{grammarlens}

\subsection*{Your turn}
\begin{itemize}
\raggedright
  \item \emph{Say it:} \emph{santākūkḍī}
  \item \emph{Say it:} \emph{santākūkḍī}, once more
  \item \emph{Say it:} \emph{bāḷko bagīcāmā̃ santākūkḍī ramtā̃ hatā̃.} — The children were playing hide-and-seek in the garden.
  \item \emph{Recall:} say the Gujarati for a spinning top, then the Gujarati for a cowrie, then say \emph{santākūkḍī} again
\end{itemize}

\begin{tcolorbox}[breakable,colback=teal!4,colframe=teal!35!black,title={Before you move on}]
What does \gu{સંતાકૂકડી} mean? (\textquotedblleft{}Hide-and-seek\textquotedblright{}.) Say it once more.
\end{tcolorbox}
\section[\texorpdfstring{\gu{કોયડો} (koyḍo)}{કોયડો (koyḍo)}]{\gu{કોયડો} (koyḍo) — a puzzle}

\label{lesson:GU-C524-koydo}



\begin{quote}
Before the new one: say the Gujarati for a cowrie, then the Gujarati for hide-and-seek.
\end{quote}

\subsection*{You'll want to know: \gu{કોયડો}}
\textbf{\gu{કોયડો}} — \emph{koyḍo} — \textquotedblleft{}a puzzle\textquotedblright{}.

\begin{quote}
\gu{અરુણે} \gu{ગણિતનો} \gu{એક} \gu{અઘરો} \gu{કોયડો} \gu{ઉકેલ્યો}.
\end{quote}

\begin{quote}
\emph{Aruṇe gaṇitno ek aghro koyḍo ukelyo.}
\end{quote}

\begin{quote}
Arun solved a hard maths puzzle.
\end{quote}

A masculine noun.

\begin{cousinweb}
\textbf{\gu{કોયડો}} (\emph{koyḍo}) is something to be solved, \textbf{\gu{ઉકેલવું}}. CDIAL files the Gujarati word, \textquotedblleft{}a puzzle, a riddle\textquotedblright{}, under Sanskrit \emph{kuha}, \textquotedblleft{}a cheat\textquotedblright{} (CDIAL 3385); how a trick became a puzzle is not explained there.
\end{cousinweb}

\begin{grammarlens}[title={What you've built}]
One more word for describing what you lived and felt.
\end{grammarlens}

\subsection*{Your turn}
\begin{itemize}
\raggedright
  \item \emph{Say it:} \emph{koyḍo}
  \item \emph{Say it:} \emph{koyḍo}, once more
  \item \emph{Say it:} \emph{Aruṇe gaṇitno ek aghro koyḍo ukelyo.} — Arun solved a hard maths puzzle.
  \item \emph{Recall:} say the Gujarati for a cowrie, then the Gujarati for hide-and-seek, then say \emph{koyḍo} again
\end{itemize}

\begin{tcolorbox}[breakable,colback=teal!4,colframe=teal!35!black,title={Before you move on}]
What does \gu{કોયડો} mean? (\textquotedblleft{}A puzzle\textquotedblright{}.) Say it once more.
\end{tcolorbox}
\section[(dialogue)]{Looking back: from Wires and Switches to Games and Collections}

\label{lesson:GU-R524-review-517-524}



\begin{quote}
Say the Gujarati for an electric wire and a plug.
\end{quote}

\subsection*{Your turn}
Say these aloud:

\begin{itemize}
\raggedright
  \item an electric wire, a plug, a switch, a light bulb, to push in (firmly)
  \item waste paper, a cobweb, a long broom, mould, to stick (something) on
  \item a charcoal stove, a tongue of flame, cooking heat, steam, to let cool
  \item to temper, to sauté, a pinch, to sprinkle, to parch
  \item a curd starter, to set curd, to churn, to boil over, to whisk
  \item a mortar, to pound, to grate, to soak, to mill grain
  \item savoury snacks, khaman, undhiyu, to be to one's taste, to dunk
  \item a collection, a spinning top, a cowrie, hide-and-seek, a puzzle
\end{itemize}

\begin{tcolorbox}[breakable,colback=teal!4,colframe=teal!35!black,title={Before you move on}]
An electric wire? (\textbf{\gu{વાયર}}, \emph{vāyar}.) A plug? (\textbf{\gu{પ્લગ}}, \emph{plag}.) A switch? (\textbf{\gu{સ્વિચ}}, \emph{svic}.) A light bulb? (\textbf{\gu{બલ્બ}}, \emph{balb}.) To push in (firmly)? (\textbf{\gu{ખોસવું}}, \emph{khosvũ}.) Waste paper? (\textbf{\gu{પસ્તી}}, \emph{pastī}.) A cobweb? (\textbf{\gu{જાળું}}, \emph{jāḷũ}.) A long broom? (\textbf{\gu{ઝાડુ}}, \emph{jhāḍu}.) Mould? (\textbf{\gu{ફૂગ}}, \emph{phūg}.) To stick (something) on? (\textbf{\gu{ચોંટાડવું}}, \emph{cõṭāḍvũ}.) A charcoal stove? (\textbf{\gu{સગડી}}, \emph{sagḍī}.) A tongue of flame? (\textbf{\gu{ઝાળ}}, \emph{jhāḷ}.) Cooking heat? (\textbf{\gu{આંચ}}, \emph{ā̃c}.) Steam? (\textbf{\gu{વરાળ}}, \emph{varāḷ}.) To let cool? (\textbf{\gu{ઠારવું}}, \emph{ṭhārvũ}.) To temper? (\textbf{\gu{વઘારવું}}, \emph{vaghārvũ}.) To sauté? (\textbf{\gu{સાંતળવું}}, \emph{sā̃taḷvũ}.) A pinch? (\textbf{\gu{ચપટી}}, \emph{capṭī}.) To sprinkle? (\textbf{\gu{છાંટવું}}, \emph{chā̃ṭvũ}.) To parch? (\textbf{\gu{ભૂંજવું}}, \emph{bhū̃jvũ}.) A curd starter? (\textbf{\gu{મેળવણ}}, \emph{meḷvaṇ}.) To set curd? (\textbf{\gu{જમાવવું}}, \emph{jamāvvũ}.) To churn? (\textbf{\gu{વલોવવું}}, \emph{valovvũ}.) To boil over? (\textbf{\gu{ઊભરાવું}}, \emph{ūbhrāvũ}.) To whisk? (\textbf{\gu{ફીણવું}}, \emph{phīṇvũ}.) A mortar? (\textbf{\gu{ખાયણી}}, \emph{khāyṇī}.) To pound? (\textbf{\gu{ખાંડવું}}, \emph{khā̃ḍvũ}.) To grate? (\textbf{\gu{છીણવું}}, \emph{chīṇvũ}.) To soak? (\textbf{\gu{પલાળવું}}, \emph{palāḷvũ}.) To mill grain? (\textbf{\gu{દળવું}}, \emph{daḷvũ}.) Savoury snacks? (\textbf{\gu{ફરસાણ}}, \emph{pharsāṇ}.) Khaman? (\textbf{\gu{ખમણ}}, \emph{khamaṇ}.) Undhiyu? (\textbf{\gu{ઊંધિયું}}, \emph{ū̃dhiyũ}.) To be to one's taste? (\textbf{\gu{ભાવવું}}, \emph{bhāvvũ}.) To dunk? (\textbf{\gu{બોળવું}}, \emph{boḷvũ}.) A collection? (\textbf{\gu{સંગ્રહ}}, \emph{sangrah}.) A spinning top? (\textbf{\gu{ભમરડો}}, \emph{bhamarḍo}.) A cowrie? (\textbf{\gu{કોડી}}, \emph{koḍī}.) Hide-and-seek? (\textbf{\gu{સંતાકૂકડી}}, \emph{santākūkḍī}.) A puzzle? (\textbf{\gu{કોયડો}}, \emph{koyḍo}.)
\end{tcolorbox}
